% Progress report for grant work packages 10 and 11 (calendar year 2026).
% Every figure quoted here is traceable to a committed result file under reports/
% in the project repository; the provenance column names that file. Where a work
% package objective has NOT been met, it is stated as such rather than softened --
% these are milestone deliverables and an inflated status here is expensive later.
\documentclass[11pt,a4paper]{article}
\usepackage[T1]{fontenc}
\usepackage[utf8]{inputenc}
\usepackage{mathptmx}
\usepackage[margin=2.2cm]{geometry}
\usepackage{microtype}
\usepackage{booktabs}
\usepackage{tabularx}
\usepackage{siunitx}
\sisetup{detect-all, per-mode=symbol, range-phrase={--}, range-units=single}
\usepackage{enumitem}
\usepackage{caption}
\usepackage{xcolor}
\usepackage[hidelinks]{hyperref}
\usepackage{fancyhdr}

\pagestyle{fancy}
\fancyhf{}
\lhead{\small Progress report: work packages 10 and 11}
\rhead{\small \thepage}
\renewcommand{\headrulewidth}{0.4pt}

\setlength{\parindent}{0pt}
\setlength{\parskip}{6pt}
\newcommand{\ms}{\ensuremath{m_s}}
\newcommand{\RMSE}{\ensuremath{\mathrm{RMSE}}}
\newcommand{\Czon}{\ensuremath{C_{\mathrm{zon}}}}
\newcommand{\done}{\textcolor{black!70}{\textbf{Delivered}}}
\newcommand{\partialx}{\textcolor{black!70}{\textbf{Partly delivered}}}
\newcommand{\notdone}{\textcolor{black!70}{\textbf{Not delivered}}}

\begin{document}

\begin{center}
{\Large\bfseries Data-Driven Results for Work Packages 10 and 11}\\[4pt]
{\large Digital twins per building with edge integration; feedback control with
continuous learning for self-optimizing buildings}\\[8pt]
{\small Reporting period: January--December 2026 \quad$\vert$\quad Prepared 12 August 2026}
\end{center}

\vspace{4pt}
\hrule
\vspace{10pt}

\section*{Scope and evidential basis}

This report summarizes the measured results obtained under work packages~10 and~11. All
experiments were run on the BOPTEST framework, an open Modelica-based benchmark for building
control, using four single-zone test cases: the air-based source case \texttt{bestest\_air} and
three hydronic cases spanning residential and commercial scale. Every number below is traceable
to a committed result file; the provenance column of each table names it.

Two objectives within these work packages were not reached in the reporting period. They are
reported as such under \emph{Objectives not reached} below rather than presented as partial successes, because
both are prerequisites for the 2027 activities and the schedule depends on an accurate status.

A peer-reviewed output covering the core of both work packages was submitted to
\emph{Energies} (MDPI) on 9~August 2026, manuscript \texttt{energies-4523055}, and is under
editorial assessment.

\section*{Work package 10: digital twins per building and edge integration}

\subsection*{10.1 Twin architectures and predictive validity}

Two surrogate architectures were developed for the source building and validated against
multi-step BOPTEST rollouts rather than one-step residuals, since a twin used for control is
evaluated on trajectories, not on single transitions.

\begin{table}[h]
\centering\small
\begin{tabularx}{\linewidth}{@{}lXr@{}}
\toprule
Twin & Architecture & 24\,h rollout \RMSE\ (\si{\celsius}) \\
\midrule
BB (hourly) & compact black-box, 8482 parameters, multi-horizon rollout objective & 1.557 \\
BB (15-min) & same architecture, matched \SI{900}{\second} corpus & 0.876 \\
GB (raw) & grey-box RC balance with neural residual heat-flow head, uncalibrated & 1.466 \\
GB (calibrated) & same, after Stage A/B/C inverse identification & \textbf{0.644} \\
\bottomrule
\end{tabularx}
\caption*{\small\RMSE{} is the multi-step rollout root-mean-square error of zone temperature over a
24-hour horizon. Provenance for all four rows:
\texttt{block1\_corpus\_matched\_comparison.csv}.}
\end{table}

The improvement from the hourly black-box baseline to the calibrated twin is
\SI{0.913}{\celsius}. A corpus-matched ablation attributes \SI{0.681}{\celsius} (74.6\%) to
corpus temporal resolution and \SI{0.232}{\celsius} (25.4\%) to the staged calibration itself. An
alternative decomposition through the uncalibrated grey-box backbone assigns 90.1\% to
calibration; the two differ because corpus and architecture interact non-additively, so the
calibration contribution is reported as a range rather than a single number.

\subsection*{10.2 Physical identifiability of the twin}

The grey-box twin carries one physically interpretable parameter, the zone thermal capacitance
\Czon, recovered by staged inverse identification from telemetry.

\begin{itemize}[leftmargin=1.2em,itemsep=2pt]
\item Identified value \Czon\ $=\SI{4.413e5}{\joule\per\kelvin}$, a $+5.1\%$ update over the
physical prior and within one standard error of the Laplace/Fisher interval.
\item Physically equivalent to approximately \SI{439}{\kilogram} of air, the expected order of
magnitude for a \SI{48}{\meter\squared} zone.
\item Directional validity: sweeping the supply-temperature command at fixed state gives the
physically correct response sign in 100\% of 400 sampled states.
\item Temperature-alignment stage reduced alignment \RMSE{} from \SI{0.374}{\celsius} to
\SI{0.232}{\celsius}, a 38.0\% improvement
(\texttt{block1\_surrogate\_final\_metrics.csv}).
\end{itemize}

\subsection*{10.3 Twins for each building: recalibration across the building family}

The objective of a twin \emph{per building} was tested by porting the identification pipeline to
three further test cases that differ in heat source, distribution path and actuator interface.
The pipeline transferred to all three.

\begin{table}[h]
\centering\small
\setlength{\tabcolsep}{4pt}
\begin{tabular}{@{}lrrrrrr@{}}
\toprule
& \multicolumn{3}{c}{Zone temperature} & \multicolumn{2}{c}{HVAC power} & \\
\cmidrule(lr){2-4}\cmidrule(lr){5-6}
Building & raw \RMSE & calibrated & gain & raw MAE & calibrated & $\rho_C$ \\
& (\si{\celsius}) & (\si{\celsius}) & (\%) & (\si{\watt}) & (\si{\watt}) & \\
\midrule
Residential, heat pump (\SI{192}{\meter\squared}) & 1.421 & 0.565 & 60.2 & 2921 & 1767 & 1.892 \\
Residential, boiler (\SI{48}{\meter\squared}) & 2.666 & 0.335 & 87.4 & 784 & 85 & 1.954 \\
Commercial, district heating (\SI{8500}{\meter\squared}) & 1.952 & 0.238 & 87.8 & --- & --- & 1.909 \\
\bottomrule
\end{tabular}
\caption*{\small $\rho_C$ is the re-identified zone capacitance as a ratio to the source building.
Provenance: \texttt{block3\_hydronic\_family\_n2\_summary.csv},
\texttt{block3\_transfer\_matrix.csv}.}
\end{table}

Two results are worth isolating. First, recalibration reduced rollout error by
\SIrange{60.2}{87.8}{\percent} on every target, so the identification pipeline is a portable
asset rather than a one-building artifact. Second, the re-identified capacitance clusters at
$\rho_C = 1.918 \pm 0.032$ across an order-of-magnitude change in floor area
(\SIrange{48}{8500}{\meter\squared}). This is reported as an effective single-zone RC signature
of the hydronic family and not as a claimed physical invariant: it rests on three cases.

\subsection*{10.4 Runtime characteristics relevant to edge deployment}

Throughput was benchmarked against the live BOPTEST runtime under an identical
\SI{900}{\second} control period, on a single CPU thread with no GPU.

\begin{table}[h]
\centering\small
\begin{tabular}{@{}lrrr@{}}
\toprule
Backend & Env.\ steps/s & Mean step (\si{\milli\second}) & Speed-up \\
\midrule
BOPTEST RTE (HTTP) & 21.0 & 15.33 & $1\times$ \\
BB surrogate & 4626 & 0.19 & $220.2\times$ \\
GB calibrated surrogate & --- & --- & $114.2\times$ \\
Role-separated hybrid & --- & --- & $85.0\times$ \\
\bottomrule
\end{tabular}
\caption*{\small Provenance: \texttt{speed\_benchmark\_table.csv}. Device \texttt{cpu},
\texttt{torch\_num\_threads = 1}. At hybrid throughput a $5\times10^{6}$-step policy-optimization
run costs about 47 minutes of environment time against roughly 66 hours on the live runtime.}
\end{table}

These figures characterize the computational footprint of the twin: an 8482-parameter model
executing a control step in \SI{0.19}{\milli\second} on one commodity CPU thread. That is the
quantitative basis on which edge deployment becomes feasible. It is not itself an edge
deployment; see \emph{Objectives not reached} below.

\section*{Work package 11: feedback control with continuous learning}

\subsection*{11.1 Closed-loop control deployed on the live runtime}

Feedback control was deployed against the live BOPTEST Run-Time Environment, a Dockerized
Modelica emulator with embedded baseline loops exposed over a REST API. At each
\SI{900}{\second} step the agent reads measurements and disturbance forecasts, and overwrites a
single control input, the supply-air temperature setpoint, leaving the remaining embedded loops
at baseline. Controllers are scored with the duration-and-severity maintenance score
$\ms = r_{\mathrm{time}} + r_{\mathrm{sev}}$, where $\ms > 1$ marks an unusable controller.

\begin{table}[h]
\centering\small
\begin{tabular}{@{}lcccc@{}}
\toprule
& \multicolumn{2}{c}{\ms\ (seed 42 $\pm$ seed s.d.)} & \multicolumn{2}{c}{Comfort violation (\%)} \\
\cmidrule(lr){2-3}\cmidrule(lr){4-5}
Training backend & peak & typical & peak & typical \\
\midrule
BB hourly & $0.073 \pm 0.019$ & $0.095 \pm 0.047$ & 1.5 & 4.4 \\
BB matched-resolution & $1.142 \pm 0.000$ & $1.211 \pm 0.000$ & 85.6 & 91.4 \\
GB calibrated & $1.046 \pm 0.022$ & $1.102 \pm 0.026$ & 77.1 & 82.4 \\
Role-separated hybrid & $\mathbf{0.087 \pm 0.026}$ & $\mathbf{0.041 \pm 0.024}$ & 4.7 & 2.4 \\
\midrule
BOPTEST built-in PI (12-month) & \multicolumn{2}{c}{0.910} & \multicolumn{2}{c}{63.6} \\
\bottomrule
\end{tabular}
\caption*{\small Dispersions are standard deviations over seeds $\{42,43,44\}$. Provenance:
\texttt{block2\_fidelity\_utility\_scatter.csv}, \texttt{block2\_thermostatic\_seed\_band.csv}.}
\end{table}

The principal scientific result of the period sits in this table. The most accurate twin
(\SI{0.644}{\celsius}) trains a controller that fails on every seed, while the least accurate
(\SI{1.557}{\celsius}) trains a usable one; making the black-box twin more accurate by refining
its resolution moves it from the usable regime into the failing one. Predictive accuracy is
therefore not a valid criterion for selecting a twin as a training environment. A rescaling
control that holds scale-free response-surface roughness fixed while varying only the integration
step reproduces the collapse, isolating the per-step increment magnitude as the operative
property. Measured relative roughness: 0.169 for the usable hourly twin against $9.4\times$ and
$7.9\times$ that for the two failing backends
(\texttt{block2\_mechanism\_surface\_sharpness.csv}).

\subsection*{11.2 Role separation as the deployed control architecture}

The architecture that resolves the conflict assigns the two conflicting requirements to two
models: the black-box twin supplies the rollout dynamics, and the frozen calibrated twin acts
only as a per-step model-disagreement reward censor, never entering the rollout or the policy
loss. Averaged over three seeds this backend leads the best single-model alternative on both
evaluation windows (0.060 against 0.072 on the peak window, 0.014 against 0.046 on the typical
window), at $85\times$ live throughput. Mean twin disagreement across hybrid traces is
\SI{0.92}{\celsius} and \SI{682}{\watt} on the peak window and \SI{1.01}{\celsius} and
\SI{735}{\watt} on the typical window (\texttt{hybrid\_disagreement\_summary.csv}).

\subsection*{11.3 Controller families and the observation interface}

\begin{itemize}[leftmargin=1.2em,itemsep=3pt]
\item \textbf{Hierarchical DRL.} A seasonal hierarchy routing to winter and summer setpoint
specialists degrades as the disagreement censor is strengthened: peak comfort violation rises
from $6.3 \pm 1.7$\% at $\lambda_T = 0$ to $24.4 \pm 3.8$\% at $\lambda_T = 0.10$, the setting
optimal for the thermostatic family. The sweep is replicated over seeds $\{42,43,44\}$ at all
four settings: the endpoint gap is 10.0 (peak) and 5.8 (typical) pooled seed standard
deviations, although the $0.05 \to 0.10$ increment sits within one pooled s.d.\ so the two
strongest weights are not separable. The regularization strength is controller-family specific
and cannot be quoted without naming its family
(\texttt{block2\_hdrl\_lambda\_sweep\_seed\_band.csv}).

\item \textbf{Multi-objective RL and the observation interface.} Widening the observation from a
5-dimensional instantaneous vector to a 17-dimensional forecast-augmented one changes the
controller from unusable to usable: \RMSE{} falls from \SI{4.96}{\celsius} to
\SI{0.72}{\celsius}, comfort violation from 74.5\% to 4.9\%, and \ms{} from 1.046 to 0.099. The
binding constraint is observation geometry, not the reward scalarization
(\texttt{block2\_morl\_comparison\_summary.csv}).

\item \textbf{Multi-objective stability.} Over five fixed seeds the neutral preference point
attains $\ms = 0.187 \pm 0.078$ (coefficient of variation 0.42, range 0.103 to 0.310). The mean
is well inside the usable band but the spread fails a pre-specified stability test, so the
preference-conditioned controller is not yet deployment-stable
(\texttt{morl\_canonical\_seedfix\_yearly\_summary.csv}).
\end{itemize}

\subsection*{11.4 Evidence bearing on the continuous-learning objective}

Three results from the period speak directly to whether a frozen policy can be deployed across a
building stock, which is the premise the continuous-learning objective is meant to relax.

\begin{enumerate}[leftmargin=1.4em,itemsep=3pt]
\item \textbf{Frozen-policy transfer fails on two of three target buildings.} Both residential
cases fall short of the pre-specified $1.25\times$-PI comfort threshold ($\ms = 0.665$ against a
threshold of 0.579; $\ms = 0.976$ against 0.938) while saving 7.3\% and 5.8\% energy
respectively. The commercial case passes the comfort threshold ($\ms = 0.431$ against 0.785) but
consumes 35.3\% more energy than the PI baseline, so it is a threshold pass and not a
deployment-ready result.

\item \textbf{The transferable asset is the pipeline, not the policy.} Against the same three
targets the identification pipeline succeeded in every case (Section~10.3). Transfer is therefore
component-level: identified physics ports across buildings, the frozen control law does not.

\item \textbf{Warm-starting from the physical twin is counterproductive.} Pretraining on the
calibrated twin and then fine-tuning on the live runtime raises \ms{} by a factor of two to three
on both windows against an otherwise identical from-scratch fine-tune. Adaptation strategies for
the continuous-learning objective should not assume that twin-pretrained weights are a useful
starting point.
\end{enumerate}

Taken together these quantify the need that work package~11 addresses: local adaptation is
required rather than optional, and the adaptation must act on the identified model rather than on
transferred policy weights.

\section*{Objectives not reached in the reporting period}

\textbf{Edge computing infrastructure integration (work package 10).} \notdone. All experiments
ran on a laptop workstation (6-core CPU, \SI{32}{\giga\byte} RAM) with the emulator in a Docker
container. No deployment to edge hardware was performed and no edge orchestration was
integrated. What the period delivered instead is the quantitative precondition: a twin of 8482
parameters executing a control step in \SI{0.19}{\milli\second} on a single CPU thread, which
establishes that the computational envelope is compatible with edge-class hardware. Integration
remains to be carried out.

\textbf{Continuous online learning (work package 11).} \partialx. Feedback control is deployed
and evaluated in closed loop, and offline adaptation stages exist within the multi-objective
pipeline, including a live-runtime fine-tuning stage. What has not been deployed is a continuous
online learning loop in which a controller adapts during operation. The transfer results of
Section~11.4 establish the requirement quantitatively and identify the correct target for
adaptation, but the adaptation mechanism itself is not yet implemented.

\section*{Summary of status against the stated deliverables}

\begin{table}[h]
\centering\small
\begin{tabularx}{\linewidth}{@{}p{0.30\linewidth}p{0.14\linewidth}X@{}}
\toprule
Deliverable & Status & Evidence \\
\midrule
Digital twins developed for each building & \done & Two architectures validated on the source case; identification pipeline ported to three further buildings with \SIrange{60.2}{87.8}{\percent} rollout-error reduction on each \\
Integration with edge computing infrastructure & \notdone & Computational envelope established (\SI{0.19}{\milli\second} per step, single CPU thread, 8482 parameters); no edge deployment performed \\
Feedback control strategies deployed & \done & Closed-loop control on the live runtime across three controller families; best configuration reaches $\ms = 0.041$ with comfort violation below 5\% at $85\times$ live throughput \\
Continuous learning for self-optimization & \partialx & Offline and live fine-tuning stages implemented; transfer study quantifies the need for local adaptation; no online learning loop in operation \\
\bottomrule
\end{tabularx}
\end{table}

\section*{Recommended priorities for the next period}

\begin{enumerate}[leftmargin=1.4em,itemsep=2pt]
\item Port the black-box twin to edge-class hardware and measure step latency, memory footprint
and thermal behaviour under sustained operation. The single-thread CPU benchmark suggests this is
tractable and should be confirmed empirically rather than assumed.
\item Implement local adaptation acting on the identified model, since Section~11.4 shows that
policy-weight transfer is the wrong mechanism and model recalibration is the one that ports.
\item Extend seed replication to the remaining controller family. The hierarchical
censor-weight sweep is now replicated over three seeds at all four settings, closing what was the
weakest link in the controller-family findings; the preference-conditioned MORL controller still
fails its pre-specified stability test and is the next target.
\item Add a stronger control baseline. Present comparisons are against the emulator's built-in
PI, which is a weak reference; a tuned rule-based or model-predictive baseline on the same twin
would put the reported gains on firmer ground.
\end{enumerate}

\vspace{6pt}
\hrule
\vspace{4pt}
{\footnotesize All quantities in this report are reproducible from committed result files in the
project repository. Figures and tables of the submitted manuscript are linked to their source
artifacts through provenance maps in the Supplementary Materials.}

\end{document}
